Un tableau pointilliste est formé de petits points de couleur dont les centres sont séparés de \qty{\v{s}}{mm}. Un visiteur a une pupille de \qty{\v{Dp}}{mm} de diamètre. On suppose que la résolution de son œil n'est limitée que par la diffraction.

\begin{enumerate}[label=\alph*)]
    \item Déterminez l'angle limite de résolution de l'œil pour une lumière de \qty{550}{nm}.
    \item À partir de quelle distance du tableau le visiteur ne peut-il plus distinguer deux points voisins ?
    \item Refaites le calcul de b) pour des points bleus (\qty{450}{nm}) et pour des points rouges (\qty{650}{nm}). Lesquels se confondent en premier quand le visiteur recule ?
\end{enumerate}
